% =====================================================================
% Relação entre os trabalhos correlatos e este trabalho
% Para o fim do capítulo de trabalhos correlatos.
% Pacotes: tabularx, booktabs (\usepackage{tabularx,booktabs}).
% Dados de cada trabalho conferidos nas fontes em 2026-10-19
% (ver docs/tcc/correlatos_notas.md).
% =====================================================================

\section{Comparação entre os trabalhos correlatos}

O Quadro~\ref{qua:correlatos} resume os trabalhos discutidos neste capítulo
segundo cinco aspectos: a fonte de informação utilizada, o tipo de texto
produzido e o seu público, as técnicas empregadas, a forma de avaliação e
o ponto que permanece em aberto em relação à proposta deste trabalho.

\begin{quadro}[htb]
\caption{Comparação entre os trabalhos correlatos e este trabalho}
\label{qua:correlatos}
\footnotesize
\begin{tabularx}{\textwidth}{@{}>{\raggedright\arraybackslash}p{2.2cm}
  >{\raggedright\arraybackslash}X >{\raggedright\arraybackslash}X
  >{\raggedright\arraybackslash}X >{\raggedright\arraybackslash}X@{}}
\toprule
\textbf{Trabalho} & \textbf{Fonte e saída} & \textbf{Técnicas} &
\textbf{Avaliação} & \textbf{O que fica em aberto} \\
\midrule
\citeonline{yang2025} (DocAgent) &
Código-fonte Python; documentação de funções, métodos e classes, voltada a desenvolvedores &
Sistema multiagente (\textit{Reader}, \textit{Searcher}, \textit{Writer}, \textit{Verifier} e \textit{Orchestrator}) com processamento topológico do código &
Completude, utilidade (LLM como juiz) e veracidade, em 366 componentes de 9 repositórios; sem avaliação humana &
Documentação técnica de código; não trata de manuais para o usuário final \\
\midrule
\citeonline{jakupov2026} (DocAider) &
Repositórios GitHub; documentação em Markdown do código, atualizada a cada \textit{pull request} &
Múltiplos agentes com GPT-4, acionados por GitHub Actions &
Relato de projeto, sem métricas publicadas na fonte consultada &
Manutenção de documentação técnica; o público continua sendo o desenvolvedor \\
\midrule
\citeonline{medeiros2024} &
Manual automotivo (Volkswagen Taos 2023) já existente; respostas a perguntas do usuário &
RAG com GPT-4 Turbo, LangChain e banco vetorial Milvus; entrada por texto, áudio e imagem &
Estudo de caso com 2 perguntas; ROUGE-L, METEOR, BERTScore e tempo de resposta &
Consulta um manual pronto; não gera o manual nem parte de artefatos de desenvolvimento \\
\midrule
\citeonline{perkusich2026} &
Questionário com 70 profissionais brasileiros &
Uso de LLMs na gestão do Scrum &
Pesquisa de opinião: 78\% relatam ganho de produtividade; 81\%, saídas ``quase corretas''; 59\%, alucinações &
Mostra o risco de erro das LLMs, mas não mede a fidelidade de um texto gerado \\
\midrule
\citeonline{gurgul2026} &
Revisão da literatura e questionário com 65 desenvolvedores &
IA generativa no ciclo de vida do software &
Mais de 70\% relatam que a IA ao menos reduziu à metade o tempo de tarefas repetitivas e de documentação &
Percepção de tempo, sem comparação controlada entre métodos de geração \\
\midrule
\textbf{Este trabalho} &
32 artefatos do Jira (épicos, histórias, \textit{bugs} e melhorias); manual do usuário final &
Redação manual, \textit{templating} (Jinja2), RAG (Cohere Embed e Command~R, Chroma) e híbrido (\textit{template} + LLM) &
Tempo e \textit{tokens} em 10 execuções, fidelidade frase a frase, cobertura, legibilidade (Flesch), BERTScore, variação, LLM como juiz e testes estatísticos &
--- \\
\bottomrule
\end{tabularx}
\legend{Fonte: elaborado pela autora.}
\end{quadro}

Os trabalhos analisados mostram que a aplicação de LLMs à documentação de
software já apresenta resultados expressivos, mas concentrada na
documentação técnica. O DocAgent \cite{yang2025} gera a documentação de
funções, métodos e classes a partir do código-fonte e alcança 95,74\% de
veracidade na melhor configuração, enquanto o DocAider \cite{jakupov2026}
automatiza a atualização da documentação de repositórios a cada
\textit{pull request}. Em ambos, a entrada é o código e o leitor é o
desenvolvedor. Nenhum deles produz instruções para quem usa o sistema sem
conhecimento técnico, que é o público do manual de usuário tratado neste
trabalho.

\citeonline{medeiros2024} aproximam-se desse público ao aplicar RAG a
manuais automotivos, mas o problema que resolvem é outro: o manual já
existe, e a LLM ajuda o usuário a encontrar a informação nele. Neste
trabalho, o manual ainda não existe. A base de conhecimento do RAG é
formada pelos artefatos de desenvolvimento registrados na plataforma de
DevOps \cite{Montgomery2024_DevOps}, e o resultado esperado é o próprio
manual. Além disso, o estudo de caso de \citeonline{medeiros2024} avalia
duas perguntas por métricas de sobreposição com respostas de referência,
sem verificar frase a frase se o texto gerado é sustentado pela fonte.

As pesquisas de \citeonline{perkusich2026} e \citeonline{gurgul2026}
completam o quadro pelo lado da prática. A segunda indica ganho de tempo
percebido em tarefas de documentação, e a primeira registra que saídas
``quase corretas'' e alucinações estão entre as principais barreiras à
adoção de LLMs. Ambas, porém, baseiam-se na opinião de profissionais. Não
comparam, sob as mesmas condições, quanto tempo cada método leva nem
quanto do texto gerado é de fato sustentado pelas fontes.

Este trabalho se posiciona nessas lacunas de três formas. Primeiro, muda a
fonte e o público: parte de artefatos de gestão ágil (épicos, histórias de
usuário, \textit{bugs} e melhorias) para produzir documentação destinada
ao usuário final. Segundo, em vez de propor uma única ferramenta, compara
quatro métodos sobre a mesma entrada: a redação manual, o
\textit{templating}, o RAG e a combinação de \textit{template} com LLM.
Isso permite separar o que cada técnica contribui para a estrutura e para
o texto. Terceiro, mede de forma direta os dois riscos apontados pela
literatura: o custo de tempo, em execuções repetidas e com teste
estatístico, e a fidelidade, pela anotação de cada frase dos manuais
quanto ao apoio nos artefatos de origem, com a distinção entre erros
intrínsecos e extrínsecos de \citeonline{Maynez2020}.
